%%%%%%%%%%%%%%%%%%%%%%%%%%%%%%%%%
%
%       EXERCÍCIO
%
%%%%%%%%%%%%%%%%%%%%%%%%%%%%%%%%%

\definevalorquestao

\begin{exercicioBanco}[\valorquestao]
Uma variável aleatória discreta \(X\) tem a seguinte função de distribuição
%
\[
    F(x) = \left\{
    \begin{array}{rl}
        0, & \text{se } x < -1, \\
        0,2, & \text{se } -1 \le x < 2, \\
        0,5, & \text{se } 2 \le x < 5, \\
        0,7, & \text{se } 5 \le x < 6, \\
        0,9, & \text{se } 6 \le x < 15, \\
        1 & \text{se } x \ge 15. \\
    \end{array}
    \right.
\]
%
Determine:
\begin{enumerate}[(a)]
    \item A função de probabilidade de \(X\).
    %
    \item \(P(X \le -2)\).
    %
    \item \(P(X < 2)\).
    %
    \item \(P(3 \le X \le 12)\).
\end{enumerate}
%
\end{exercicioBanco}

